\documentclass[10pt,a4paper]{amsart}
\usepackage[margin=19mm]{geometry}
\usepackage{amsmath,amssymb,mathtools}
\usepackage[scr=rsfs,cal=euler]{mathalfa}
\usepackage{xfrac}
\usepackage[unicode,hidelinks,bookmarks=false]{hyperref}
\newcommand{\ie}{i.e.\ }
\newcommand{\cf}{cf.\ }
\newcommand{\resp}{respectively\ }
\newcommand{\Subsection}{\subsection}
\providecommand{\nobreakdash}{\nobreak\hbox{-}}
\setlength{\emergencystretch}{2em}
\multlinegap0pt
\usepackage{enumerate}
\usepackage{amssymb}
\newtheorem{proposition}{Proposition}
\newtheorem{lemma}{Lemma}
\newtheorem{theorem}{Theorem}
\newcommand{\N}{\ensuremath{\mathbb{N}}}
\renewcommand{\a}{\ensuremath{\alpha}}
\renewcommand{\b}{\ensuremath{\beta}}
\newcommand{\bbb}{\ensuremath{\mathcal{B}}}
\renewcommand{\inf}{\myinf}
\renewcommand{\l}{\ensuremath{\lambda}}
\DeclareMathOperator*{\myinf}{in\vphantom{p}f}
\def\rw{\rightarrow}
\newcommand{\ug}{\ensuremath{\mathsf{UG}}}
\newcommand{\ve}{\ensuremath{\varepsilon}}
\title{A characterisation of probabilistic metrizability for approach spaces}
\author{Eva Colebunders, Robert Lowen}
\date{Source: arXiv:2601.06269v1 (CC BY 4.0)}
\begin{document}
\maketitle

\noindent\emph{Source and licence.} A characterisation of probabilistic metrizability for approach spaces. Version arXiv:2601.06269v1 (CC BY 4.0), distributed by arXiv under the Creative Commons Attribution 4.0 International licence, http://creativecommons.org/licenses/by/4.0/, as recorded in the arXiv licence metadata for this exact version and shown on the abstract page https://arxiv.org/abs/2601.06269v1. Full licence notice: \texttt{licenses/arxiv-2601.06269-CC-BY-4.0.txt}.\par\smallskip

\noindent\textit{Editorial scope.} Counted unit: the article's theorem pmetr (source0.tex lines 660--722). The article's complete original proof of it is delivered with it (source0.tex lines 924--945, one \texttt{proof} environment of 22 lines, which the article places away from the statement). No mathematics is written, repaired or re-derived: the statement and the proof are the article's own lines, in the article's own notation, and no retained line is substituted or re-flowed.  Retained uncounted as the source-local context the proof needs: lines 379--381; lines 461--463; lines 472--474; lines 612--614; lines 618--633; lines 636--657; lines 802--804; lines 813--815; lines 902--904; lines 909--915.  Nothing was omitted from any retained span, so the package carries no omission record.  No retained line contains a citation, so the wrapper carries no bibliography and no bibliography entry.  No retained line contains a figure environment, an image-inclusion command or drawing code, so no image is delivered and the package declares no asset dependency.  Editorial formatting only, in addition: an AMS \texttt{amsart} preamble that restates the article's own declarations and macros used by the retained lines, namely the environments \texttt{proposition}, \texttt{lemma}, \texttt{theorem} on separate counters, together with the standard packages those lines need.  The retained environments print in this document as Proposition 1, Lemma 1, Theorem 1 in order of appearance, whereas the article prints the same items with the numbering of its own full document.  The complete native source of the article as distributed in the arXiv e-print for this exact version is bundled unchanged as \texttt{sources/192397/source0.tex}.
\begin{equation}\label{PrMet}
\a: X \times X \times [0,\infty] \rw [0,1]
\end{equation}

\begin{equation}\label{associated}
\bbb(x) = \{ \varphi_{\l,x} \ | \ 0 < \l \leq 1 \} \ \ 
\end{equation}

\begin{equation}\label{equivlocal}
\varphi_{\l,x} (y) < \gamma \Leftrightarrow \a(x,y,\gamma) > 1-\l.
\end{equation}

\begin{equation}\label{beta}
\b(x,y,\gamma) = \sup \{1 - \l \ | \ \varphi_{\l,x} (y) < \gamma\} = 1 - \inf  \{\l \ | \ \varphi_{\l,x} (y) < \gamma\}.
\end{equation}

\begin{proposition}\label{equivbeta}
Given an approach space $X$ having a local basis satisfying {\bf(LD)},
then for $0 \leq \gamma < \infty$ we have:
$$
\b(x,y,\gamma) > 1 - \l   \Leftrightarrow  \varphi_{\l,x} (y) < \gamma.
$$
\begin{proof}
Let  $\sup \{1 - \rho \ | \ \varphi_{\rho,x} (y) < \gamma\}  > 1 - \l.$ Then $\exists \rho < \l, \ \varphi_{\rho,x}(y) < \gamma$ and applying {\bf(LD)} we have 
$\varphi_{\l,x}(y) < \gamma.$

For the other implication,  assume $\varphi_{\l,x}(y) < \gamma.$ Then again applying {\bf(LD)}\\
 $\exists \rho < \l, 
 \  \ \varphi_{\rho,x}(y) < \gamma.$
Hence $1 - \l < 1 - \rho \leq \b(x,y,\gamma).$
\end{proof}
\end{proposition}

\begin{lemma}\label{star}
Let $([0,1], *)$ be a continuous t-norm. For all $a,b,d \in ]0,1],$ the following are equivalent:
\begin{enumerate}
\item $d \geq a * b$
\item $\forall \l,\l' \in ]0,1], $\  if $a > 1-\l, \ b > 1-\l'$ then $d \geq (1-\l') * (1-\l)$
\item $\forall \rho  \in ]0,1],$ \ if $a * b > 1-\rho$ then $d \geq 1-\rho.$
\end{enumerate}
\begin{proof}
That (1) and (3) are equivalent and that (1) implies (2) are straightforward.
We prove that (2) implies (3). 

Let $ \rho  \in ]0,1],$ and assume that $a*b > 1-\rho.$
Choose strictly increasing sequences $(1-\l_n)_n,$ in $]0,1]$ converging to $a$ and $(1-\l'_n)_n,$ converging to $b.$ 
By the continuity of $*$ on $[0,1]\times [0,1]$ the sequence $((1-\l_n) * (1-\l'_n))_n$ converges to $a * b$ with $(1-\l_n) * (1-\l'_n) \leq a*b$ for all $n.$
Since $]1-\rho, 1]$ is an open neighborhood of $a * b$ we can fix $n \in \N$ such that $(1-\l_n) * (1-\l'_n) \in ]1-\rho, 1].$
By application of (2) we now have
$$
d \geq (1-\l_n) * (1-\l'_n) > 1 - \rho.
$$

\end{proof}
\end{lemma}

\begin{equation}\label{uassociated}
\Psi = \{ d_\l \ | \ 0 < \l \leq 1 \}, \ \ 
\end{equation}

\begin{equation}\label{equivug}
d_\l (x,y) < \gamma \Leftrightarrow \a(x,y,\gamma) > 1-\l.
\end{equation}

\begin{equation}\label{ubeta}
\b(x,y,\gamma) = \sup \{1 - \l \ | \ d_\l(x, y) < \gamma\} = 1 - \inf  \{\l \ | \ d_\l (x, y) < \gamma\}.
\end{equation}

\begin{proposition}\label{uequivbeta}
Given $X$ having a uniform basis satisfying {\bf(UD)},
for $0 \leq \gamma < \infty$ we have:
$$
\b(x,y,\gamma) > 1 - \l   \Leftrightarrow  d_\l(x, y) < \gamma.
$$
\end{proposition}

\begin{theorem} [\bf Part II of Theorem 3.8]\label{pmetr}
If $X$ is an approach space having a local basis in $x$
$$
\bbb(x) = \{ \varphi_{\l,x} \ | \ 0  < \l \leq 1 \},
$$
satisfying {\bf(LS)}, {\bf(LD)}, {\bf(LT)} for some continuous t-norm $*$ and {\bf(LH)}, then there exists a probabilistic metric space $(X,\b,*)$ with underlying approach space $X$.

\begin{proof}
Suppose 
$X$ has a local basis in $x$ 
$$
\bbb(x) = \{ \varphi_{\l,x} \ | \ 0  < \l \leq 1 \},
$$
satisfying {\bf(LS), (LD), (LT)} and {\bf(LH)}.

On $X$ define $\b(x,y, -)$ as in \eqref{beta}.
We prove that $(X, \b, *) $ is a probabilistic metric space in the sense of \eqref{PrMet}.

$\b(x,y,0) =0$ as $\varphi_{\l,x} (y) \geq 0$ and $\b(x,y,\infty) =1$ by definition.

$\b(x,y,-)$ is nondecreasing as for $\gamma \leq \gamma',$  $\{ \l \ | \ \varphi_{\l,x} (y) < \gamma\} \subseteq \{ \l \ | \ \varphi_{\l,x} (y) < \gamma'\}.$


$\b(x,y,-)$ is left continuous on $]0, \infty[$. That $\sup_{s < \gamma} \b(x,y,s) \leq  \b(x,y,\gamma)$ is clear from the previous property. For the reverse inequality assume $\b(x,y,\gamma) > 1 - \l$. By \ref{equivbeta} this implies that $\varphi _{ \l,x} (y) < \gamma$  and hence $\exists s < \gamma,  \ \ \varphi _{ \l,x} (y) < s.$ This implies that $\exists s <\gamma,  \  \b(x,y,s) > 1 - \l.$ So we have that $\sup_{s < \gamma} \b(x,y,s) > 1 - \l.$
From the previous claims we can conclude that (P1) holds.

To see that (P2) holds,  observe that for $\gamma >0$  we have\\
 $$\b(x,x,\gamma) = \sup \{1 - \l \ | \ \varphi_{\l,x} (x) < \gamma\}  = 1.$$

Clearly (P3) follows immediately from {\bf(LS)}.

Next we check (P4). Suppose $x \not = y,$  from {\bf (LH)} it follows that $$\exists \l, \ \varphi_{\l,x} (y) \not = 0.$$
Choosing $\gamma$ such that $0 < \gamma < \varphi_{\l,x} (y), $ from \ref{equivbeta} we obtain $\b(x,y,\gamma) \not = 1.$

Next we check that $\b$ satisfies (P5). Let $x,y,z \in X,$ and assume $\b(y,z,\gamma') > 1 - \l'$ and $\b(x,y,\gamma) > 1 - \l.$  
For every $ 0 < \ve \leq 1 $ satisfying 
$$
(1 - \l') * ( 1 - \l) > 1-\ve,
$$
applying {\bf(LT)},  
we have
$$
\varphi_{\ve,x} (z) \leq \varphi _{\l, x} (y) + \varphi _{\l', y} (z) < \gamma' + \gamma.
$$
Hence
$$
\b(x,z, \gamma + \gamma') > 1 - \ve. 
$$
This implies
$$
\b(x,z, \gamma + \gamma') \geq  (1 - \l') * ( 1 - \l).
$$
Applying \ref{star}
we can conclude that $\b(x,z, \gamma + \gamma') \geq \b(y,z, \gamma')  * \b(x,y, \gamma). $\\

Finally we prove that the underlying approach space of $(X,\b)$ is $X$.\\
Let $\{\psi_{\l,x} \ | \ 0 < \l \leq 1 \}$ be associated with $(X,\b)$ as in \eqref{associated}. Then from \eqref{equivlocal}  and \ref{equivbeta} for $\gamma$ finite
$$
\psi_{\l,x}(y) < \gamma \Leftrightarrow \b(x,y, \gamma) > 1 - \l \Leftrightarrow \varphi_{\l,x}(y) < \gamma.
$$
Hence we can conclude that \ $\psi_{\l,x} = \varphi_{\l,x}$ for all $\l$. 
\end{proof}
\end{theorem}

\begin{proof}
Suppose
$X$ has a uniform basis 
$$
\Psi = \{ d_\l  \ | \ 0  < \l \leq 1 \}
$$
satisfying {\bf(US), (UD), (UT)} and {\bf(UH)}.

On $X$ define $\b(x,y, -)$ as in \eqref{ubeta}.
Similar to the proof in \ref{pmetr},
$(X, \b, *) $ is a probabilistic metric space.



Finally we prove that the underlying $\ug$-space of $(X,\b)$ is $X$.

Let $\{ e_\l \ | \ 0 < \l \leq 1 \}$ be associated with $(X,\b)$ as in \eqref{uassociated}. Then from \eqref{equivug}  and \ref{uequivbeta} for $\gamma<\infty$ 
$$
e_\l(x, y) < \gamma \Leftrightarrow \b(x,y, \gamma) > 1 - \l \Leftrightarrow d_\l(x, y) < \gamma.
$$
Hence we can conclude that $e_\l = d_{\l}$ for all $\l$.
\end{proof}

\end{document}
